\subsection{基于内容的引用预测}
早期工作将文献表示为词袋向量，通过相似度检索候选引用\cite{kipf2017}。这类方法实现简单，但受限于词汇鸿沟，难以识别表述不同但语义相近的工作。随着词嵌入与句向量的发展，语义匹配精度有所提升，但仍缺乏对全局引用结构的利用。

\subsection{基于图结构的引用预测}
将文献视为节点、引用关系视为边，引用预测可被形式化为链路预测问题。图卷积网络（GCN）\cite{kipf2017} 与图注意力网络（GAT）通过聚合邻居表征学习节点嵌入。然而，真实引文图是\emph{异构}的：除了引用边，还存在作者—论文、会议—论文等多种关系。异构图注意力网络（HAN）\cite{wang2019} 通过为不同元路径分配注意力权重，缓解了这一问题，但其注意力机制本身仍依赖输入节点的初始特征质量。

\subsection{预训练语言模型与中文处理}
以 BERT\cite{devlin2019} 与 ERNIE\cite{sun2020} 为代表的预训练模型，在中文自然语言理解任务上取得了突破性进展。本文采用这类模型的中文变体作为节点特征提取器，为图模型提供语义层面的初始化。
